% -----------------------------
% 執筆にあたっての留意点
% -----------------------------
% 実際にアブストラクトを執筆する際には，以下の項目を順番に1〜2行ずつ書くと良いでしょう．

% 1. 論文の目的（何を提案する，何を調査するのか）の要約
% 2. 提案手法・調査事項の具体的なアイデア
% 3. 評価実験・調査の結果，明らかになったこと
% 4. 本論文によってどんなことが可能になるか，社会や学術コミュニティにどんな波及効果があるか
%
%
% -----------------------------
% 補足（参考：how-to-xx/how-to-write-a-paper/1-学術論文の書き方.md 3.2節）
% -----------------------------
%
% ■ アブストラクトは論文全体のミニチュアです
% 読者や査読者の多くは，「アブストラクト → 図 → 結論 →（気になれば）本文」という順に流し読みします．
% つまり，多くの人にとってアブストラクトが唯一読まれる部分です．
% ここだけを読んでも，何を問題とし，何を提案し，何がわかったのかが伝わる状態を目指してください．
%
% ■ 背景や動機を長々と書いてはいけません
% 学生が最もやりがちな失敗は，背景の説明にアブストラクトの半分を使ってしまうことです．
% 背景に割いてよいのは，多くとも1文です．上の項目1〜4のうち，紙面の大半は項目2と項目3に使ってください．
% 分量の目安は次の通りです．
%
% * 背景：0〜1文
% * 問題：1文
% * 既存手法の限界：1文
% * 提案：1〜2文
% * 結果：1〜2文
%
% 和文アブストラクトは300〜500字が目安です．
% この分量の目安に沿って1〜2文ずつ書いていけば，おおむねこの範囲に収まります
% （下の例は約440字です）．
%
% ■ 結果は，可能な限り数値で書きます
% 「大幅に改善した」ではなく「従来手法と比べて◯%改善した」「検索回数が◯倍に増えた」と書くと，
% 主張が一気に強くなります．数値のないアブストラクトは，読者に「本当に効果があったのか」を判断させられません．
%
% ■ 本文の文を一語一句そのままコピーして使わないでください
% アブストラクトはイントロダクションの要約ではなく，独立した一つの文章として書きます．
% イントロダクションの第1段落をそのまま貼り付けたものは，すぐに見抜かれます．
%
% ■ その他の作法
% * 原則として，アブストラクトには文献を引用しません（引用がないと意味の通らない書き方を避けます）
% * 略語は，アブストラクト内で初めて使うときに展開します
%   （本文でも改めて展開します．アブストラクトと本文は，初出を別々に数えます）
% * 図表や数式への参照（「図1に示すように」など）も書きません
% * 書くのは最後です．本文と結論を書き終えてから，論文全体を眺めて書くのが最も速く，正確です
%   （ただし「研究を始めた時点から書き始める論文」では，現時点で書けることを書いて，何度も更新していきます）
%
%
% 以下は例です．
% 各文末の「（項目◯に対応）」は，上の項目との対応を示すための注記です．
% 自分の論文を書く際には削除してください．

本稿では，ウェブ検索中のユーザに注意深い情報探索を暗黙的に促す「クエリプライミング」を提案する（項目1に対応）．
クエリプライミングは，批判的思考を喚起し注意深い情報探索や意思決定を促進するようなキーワードを，検索ワード補完・推薦時に提示する（項目2に対応）．
クエリプライミングの有効性を検証するために，クラウドソーシングを用いたオンラインユーザ実験を行った．
被験者の情報探索ログ分析および実験のアンケート調査の結果，以下の傾向が明らかになった：
(1) クエリプライミングが実装されたウェブ検索エンジンを用いた被験者はセッション中の検索回数が増え，検索結果一覧ページを何度も見直すようになる．
(2) クエリプライミングによって，証拠を重視してウェブページを検索・閲覧する行動が促進される．
(3) クエリプライミングの効果は被験者の学歴に依存する（項目3に対応）．
本研究で得られた知見によって，注意深い情報探索を活性化させる検索インタラクションの設計に寄与することが期待される（項目4に対応）．
